
\subsubsection{4.1 Dataset}

\textbf{Waterbirds} \citep{Sagawaetal2020}: A spurious correlation benchmark constructed by superimposing CUB-200-2011 bird images onto Places backgrounds. The training split contains a 95\% spurious correlation between bird species (landbird/waterbird — the core label) and background type (land/water — the spurious label). Standard splits are used: 4,795 training, 1,199 validation, 5,794 test images. The validation split serves for probe fitting (80\%) and probe evaluation (20\%), as well as for DFR last-layer reweighting.

\textbf{CelebA} \citep{Liuetal2021}: Planned as a replication dataset with hair color (spurious) versus gender (core) as the spurious correlation pair. Not executed due to Google Drive access failure in the experimental environment. This constitutes a confirmed network restriction, not a dataset issue.

\subsubsection{4.2 Evaluation Metrics}

\textbf{Worst-group accuracy (WGA):} Accuracy on the minority group (landbirds on water backgrounds; waterbirds on land backgrounds), the standard robustness metric for spurious correlation methods.

\textbf{δ(t) metrics:} Contiguous window fraction (proportion of training epochs in the longest δ(t)>0 window), gap area A, one-sided paired t-test p-value and t-statistic across seeds.

\textbf{GDR metrics:} Mean early GDR across seeds (epochs 2, 4, 6), Wilcoxon signed-rank test p-value.

\textbf{Complexity metrics:} Uncorrected and Bonferroni-corrected (α=0.0167) p-values for each of three metrics.

*\textit{t\} metrics:** Mean and standard deviation across seeds, 95\% bootstrap CI for std (10,000 resamples).

\textbf{DFR metrics:} DFR WGA at each checkpoint, DFR improvement = DFR WGA − ERM WGA, Pearson r between improvement and epochs-past-t\*.

